% Trascrizione quiz Paola Suria: intervallo pagine del PDF.
% Ogni risposta e spiegazione e` conservata nei commenti % SOLUZIONE,
% per non alterare il parser del sito esistente.
% ATTENZIONE: verificare le voci segnalate come da verificare.

\par\noindent\hfill{\color{cyan!65!black}\scriptsize Pag. originale 12/521}\par\smallskip

% SIMULAZIONE / QUESITO 1 (pagina 12)
\begin{quiz}
Il dominio della funzione $f(x)=\log(\cosh x-4)$
\begin{enumerate}[label=\alph*),leftmargin=1.3cm,itemsep=3pt,topsep=5pt]
\item \`e $\mathbb R$
\item \`e $[0,+\infty)$
\item non \`e un intervallo
\item \`e limitato
\item \`e $[-4,4]$
\end{enumerate}
\end{quiz}
% SOLUZIONE: c
% SPIEGAZIONE: $\cosh x>4$ equivale a $|x|>\operatorname{arcosh}4$: il dominio \`e unione di due intervalli.

% SIMULAZIONE / QUESITO 2 (pagina 12)
\begin{quiz}
La funzione $f(x)=\dfrac1{4+x^2}$
\begin{enumerate}[label=\alph*),leftmargin=1.3cm,itemsep=3pt,topsep=5pt]
\item \`e invertibile su $\mathbb R$
\item \`e invertibile su $[0,+\infty)$
\item \`e invertibile su $[-2,2]$
\item non \`e invertibile su alcun intervallo
\item nessuna delle altre risposte
\end{enumerate}
\end{quiz}
% SOLUZIONE: b
% SPIEGAZIONE: Su $[0,+\infty)$ \`e strettamente decrescente, dunque iniettiva.

% SIMULAZIONE / QUESITO 3 (pagina 12)
\begin{quiz}
La disequazione in campo complesso $|z-1+i|<3$
\begin{enumerate}[label=\alph*),leftmargin=1.3cm,itemsep=3pt,topsep=5pt]
\item ha infinite soluzioni reali
\item non ha soluzioni
\item non ha soluzioni reali ma ha soluzioni complesse
\item ha un'unica soluzione
\item ha esattamente due soluzioni
\end{enumerate}
\end{quiz}
% SOLUZIONE: a
% SPIEGAZIONE: Per $z=x\in\mathbb R$, $(x-1)^2+1<9$, cio\`e $1-\sqrt8<x<1+\sqrt8$.

% SIMULAZIONE / QUESITO 4 (pagina 12)
\begin{quiz}
Il limite $\displaystyle\lim_{x\to+\infty}x^3\bigl(\log(4x^3+1)-\log(4x^3)\bigr)$ vale
\begin{enumerate}[label=\alph*),leftmargin=1.3cm,itemsep=3pt,topsep=5pt]
\item $0$
\item $+\infty$
\item $1$
\item $1/4$
\item $e^3$
\end{enumerate}
\end{quiz}
% SOLUZIONE: d
% SPIEGAZIONE: $x^3\log(1+1/(4x^3))\to1/4$ per $\log(1+t)\sim t$.

% SIMULAZIONE / QUESITO 5 (pagina 12)
\begin{quiz}
Il limite $\displaystyle\lim_{x\to+\infty}(\sqrt[4]{x^4+2x^3}-x)$ vale
\begin{enumerate}[label=\alph*),leftmargin=1.3cm,itemsep=3pt,topsep=5pt]
\item $+\infty$
\item $-\infty$
\item $0$
\item $1$
\item $1/2$
\end{enumerate}
\end{quiz}
% SOLUZIONE: e
% SPIEGAZIONE: $x[(1+2/x)^{1/4}-1]\to 1/2$.

\par\noindent\hfill{\color{cyan!65!black}\scriptsize Pag. originale 13/521}\par\smallskip

% SIMULAZIONE / QUESITO 6 (pagina 13)
\begin{quiz}
La funzione $f(x)=\log(3+2e^x)+6x$ ha come asintoto obliquo destro la retta
\begin{enumerate}[label=\alph*),leftmargin=1.3cm,itemsep=3pt,topsep=5pt]
\item $y=2x+\log3$
\item $y=7x+\log2$
\item $y=8x+\log2$
\item $y=2x$
\item $y=6x$
\end{enumerate}
\end{quiz}
% SOLUZIONE: b
% SPIEGAZIONE: $\log(3+2e^x)=x+\log(2+3e^{-x})=x+\log2+o(1)$.

% SIMULAZIONE / QUESITO 7 (pagina 13)
\begin{quiz}
Sia $a_n=n\cos(\pi n^2)$. Allora
\begin{enumerate}[label=\alph*),leftmargin=1.3cm,itemsep=3pt,topsep=5pt]
\item $a_n$ non ha limite
\item $a_n$ \`e positiva $\forall n$
\item $a_n$ \`e limitata
\item $a_n$ \`e monotona
\item $a_n$ tende sia a $+\infty$ che a $-\infty$
\end{enumerate}
\end{quiz}
% SOLUZIONE: a
% SPIEGAZIONE: $\cos(\pi n^2)=(-1)^n$, dunque $a_n=(-1)^n n$ oscilla senza limite.

% SIMULAZIONE / QUESITO 8 (pagina 13)
\begin{quiz}
Sia $f(x)=(3\sin x)(5\sin x)$. La derivata di $f(x)$ calcolata per $x=\pi/2$ vale
\begin{enumerate}[label=\alph*),leftmargin=1.3cm,itemsep=3pt,topsep=5pt]
\item $\pi$
\item $35$
\item $1$
\item $0$
\item $-\pi$
\end{enumerate}
\end{quiz}
% SOLUZIONE: d
% SPIEGAZIONE: $f(x)=15\sin^2x$, $f\prime(x)=30\sin x\cos x$, che si annulla in $\pi/2$.

% SIMULAZIONE / QUESITO 9 (pagina 13)
\begin{quiz}
Sia $f:\mathbb R\to\mathbb R$ derivabile, con $f(0)=3$ e $f\prime(0)=-4$. Allora $g(x)=|f(x)|$
\begin{enumerate}[label=\alph*),leftmargin=1.3cm,itemsep=3pt,topsep=5pt]
\item ha un punto angoloso in $x=0$
\item \`e derivabile nel punto $x=0$
\item ha un massimo relativo nel punto $x=0$
\item ha una discontinuit\`a di tipo salto nel punto $x=0$
\item ha un asintoto verticale nel punto $x=0$
\end{enumerate}
\end{quiz}
% SOLUZIONE: b
% SPIEGAZIONE: Poich\'e $f(0)>0$, vicino a zero $|f|=f$ e quindi $g\prime(0)=-4$.

% SIMULAZIONE / QUESITO 10 (pagina 13)
\begin{quiz}
Sia $f(x)=2e^x+3x$. Allora, osservando che $f(0)=2$, risulta
\begin{enumerate}[label=\alph*),leftmargin=1.3cm,itemsep=3pt,topsep=5pt]
\item $(f^{-1})\prime(2)=2$
\item $(f^{-1})\prime(2)=1/2$
\item $(f^{-1})\prime(2)=1$
\item $(f^{-1})\prime(2)=1/5$
\item $(f^{-1})\prime(2)=5$
\end{enumerate}
\end{quiz}
% SOLUZIONE: d
% SPIEGAZIONE: $(f^{-1})\prime(2)=1/f\prime(0)=1/(2+3)=1/5$.

\par\noindent\hfill{\color{cyan!65!black}\scriptsize Pag. originale 14/521}\par\smallskip

% SIMULAZIONE / QUESITO 11 (pagina 14)
\begin{quiz}
Siano $f(x)=\sqrt{1+x^4}$ e $g(x)=e^{x^4}$. Allora
\begin{enumerate}[label=\alph*),leftmargin=1.3cm,itemsep=3pt,topsep=5pt]
\item $g(x)-f(x)\le0$ in un intorno di $x=0$
\item $g(x)-f(x)=o(x^4)$ per $x\to0$
\item $g(x)-f(x)=o(x^8)$ per $x\to0$
\item $g(x)-f(x)\ge0$ in un intorno di $x=0$
\item nessuna delle altre risposte
\end{enumerate}
\end{quiz}
% SOLUZIONE: d
% SPIEGAZIONE: $e^{x^4}=1+x^4+o(x^4)$ e $\sqrt{1+x^4}=1+x^4/2+o(x^4)$; la differenza \`e positiva.

% SIMULAZIONE / QUESITO 12 (pagina 14)
\begin{quiz}
Se $f$ ha sviluppo di Taylor $f(x)=2-3(x-2)-5(x-2)^2+(x-2)^3+o((x-2)^3)$ per $x\to2$, allora in un intorno di $x=2$
\begin{enumerate}[label=\alph*),leftmargin=1.3cm,itemsep=3pt,topsep=5pt]
\item $f$ \`e negativa, crescente e concava
\item $f$ \`e negativa, crescente e convessa
\item $f$ \`e negativa, decrescente e convessa
\item $f$ \`e positiva, crescente e concava
\item $f$ \`e negativa, decrescente e convessa
\end{enumerate}
\end{quiz}
% SOLUZIONE: nessuna
% SPIEGAZIONE: $f(2)=2>0$, $f\prime(2)=-3<0$, $f\prime\prime(2)=-10<0$: positiva, decrescente e concava. Nessuna opzione coincide.

% SIMULAZIONE / QUESITO 13 (pagina 14)
\begin{quiz}
La funzione $f:\mathbb R\to\mathbb R$, $f(x)=x^5+\sinh x$, assume il valore $y=10^5+\sinh(10)-1$
\begin{enumerate}[label=\alph*),leftmargin=1.3cm,itemsep=3pt,topsep=5pt]
\item in un punto $c>10$
\item in nessun punto
\item in almeno due punti
\item in un punto $c<10$
\item nessuna delle altre risposte \`e corretta
\end{enumerate}
\end{quiz}
% SOLUZIONE: d
% SPIEGAZIONE: $f\prime(x)=5x^4+\cosh x>0$, quindi $f$ \`e strettamente crescente e il valore inferiore a $f(10)$ \`e assunto prima di $10$.

% SIMULAZIONE / QUESITO 14 (pagina 14)
\begin{quiz}
Sia $f:\mathbb R\to\mathbb R$ derivabile tale che $f(0)=f(1)=0$. Ponendo $g(x)=f(x)^3$, allora
\begin{enumerate}[label=\alph*),leftmargin=1.3cm,itemsep=3pt,topsep=5pt]
\item la derivata $g\prime(x)$ si annulla almeno tre volte
\item la derivata $g\prime(x)$ si annulla esattamente due volte
\item la derivata $g\prime(x)$ si annulla esattamente una volta
\item la derivata $g\prime(x)$ non si annulla mai
\item nessuna delle altre affermazioni \`e corretta
\end{enumerate}
\end{quiz}
% SOLUZIONE: a
% SPIEGAZIONE: $g\prime=3f^2f\prime$ si annulla in $0$, in $1$ e in un punto interno per Rolle.

% SIMULAZIONE / QUESITO 15 (pagina 14)
\begin{quiz}
Il polinomio di MacLaurin di ordine 3 della funzione $f(x)=e^{\sqrt{1+3x^3}}$ \`e
\begin{enumerate}[label=\alph*),leftmargin=1.3cm,itemsep=3pt,topsep=5pt]
\item $1+\frac32x^3$
\item $e+\frac{3e}{2}x^3$
\item $2+\frac12x^3$
\item $1+\frac e2x^3$
\item $1+\frac15x^3$
\end{enumerate}
\end{quiz}
% SOLUZIONE: b
% SPIEGAZIONE: $\sqrt{1+3x^3}=1+\frac32x^3+o(x^3)$; dunque $e^{\sqrt{1+3x^3}}=e\bigl(1+\frac32x^3+o(x^3)\bigr)$. Il polinomio di MacLaurin di ordine 3 e` $e+\frac{3e}{2}x^3$.

\par\noindent\hfill{\color{cyan!65!black}\scriptsize Pag. originale 15/521}\par\smallskip

% SIMULAZIONE / QUESITO 16 (pagina 15)
\begin{quiz}
Sia $f$ una funzione tale che $f(x)=o(x^3)$ per $x\to+\infty$. Allora possiamo sicuramente affermare che
\begin{enumerate}[label=\alph*),leftmargin=1.3cm,itemsep=3pt,topsep=5pt]
\item $\lim_{x\to+\infty}\frac{f(x)}{x^2+x^3}=+\infty$
\item $\lim_{x\to+\infty}\frac{f(x)}{x^3}=1$
\item $\lim_{x\to+\infty}\frac{f(x)}{x^3+x^2}=0$
\item $\lim_{x\to+\infty}\frac{f(x)}{x^2}=0$
\item $\lim_{x\to+\infty}\frac{f(x)}{x^3}=+\infty$
\end{enumerate}
\end{quiz}
% SOLUZIONE: c
% SPIEGAZIONE: Per definizione di $o(x^3)$, $f(x)/x^3\to0$. Poiche` $\frac{f(x)}{x^3+x^2}=\frac{f(x)/x^3}{1+1/x}$ e il denominatore tende a 1, il limite richiesto e` 0.

% SIMULAZIONE / QUESITO 17 (pagina 15)
\begin{quiz}
L\'integrale $\displaystyle\int_0^{2\pi}|\cos x|\,dx$ vale
\begin{enumerate}[label=\alph*),leftmargin=1.3cm,itemsep=3pt,topsep=5pt]
\item $2$
\item $\pi$
\item $4$
\item $2\pi$
\item $4\pi$
\end{enumerate}
\end{quiz}
% SOLUZIONE: c
% SPIEGAZIONE: Per simmetria, $4\int_0^{\pi/2}\cos x\,dx=4$.

% SIMULAZIONE / QUESITO 18 (pagina 15)
\begin{quiz}
La funzione $F(x)=\displaystyle\int_0^x t^4\cos^4(t)\,dt$
\begin{enumerate}[label=\alph*),leftmargin=1.3cm,itemsep=3pt,topsep=5pt]
\item \`e crescente su $\mathbb R$
\item \`e crescente solo su $(0,+\infty)$
\item \`e crescente solo su $(-\infty,0)$
\item \`e sempre positiva
\item non si annulla mai
\end{enumerate}
\end{quiz}
% SOLUZIONE: a
% SPIEGAZIONE: $F\prime(x)=x^4\cos^4x\ge0$, quindi $F$ \`e non decrescente su $\mathbb R$.

% SIMULAZIONE / QUESITO 19 (pagina 15)
\begin{quiz}
L\'integrale improprio $\displaystyle\int_0^1\frac{x}{1-\cos x}\,dx$
\begin{enumerate}[label=\alph*),leftmargin=1.3cm,itemsep=3pt,topsep=5pt]
\item diverge a $+\infty$
\item diverge a $-\infty$
\item \`e indeterminato
\item converge a un valore positivo
\item converge a un valore negativo
\end{enumerate}
\end{quiz}
% SOLUZIONE: a
% SPIEGAZIONE: Vicino a zero, $1-\cos x\sim x^2/2$ e l\'integranda \`e equivalente a $2/x$.

% SIMULAZIONE / QUESITO 20 (pagina 15)
\begin{quiz}
L\'integrale generale dell\'equazione differenziale $y\prime\prime-4y\prime+4y=0$ \`e
\begin{enumerate}[label=\alph*),leftmargin=1.3cm,itemsep=3pt,topsep=5pt]
\item $y=C_1x+C_2e^{2x}$
\item $y=C_1e^{-2x}+C_2e^{2x}$
\item $y=C_1xe^x+C_2e^{2x}$
\item $y=(C_1+C_2x)e^{2x}$
\item $y=C_1+C_2x$
\end{enumerate}
\end{quiz}
% SOLUZIONE: d
% SPIEGAZIONE: Il polinomio caratteristico \`e $(r-2)^2$, radice doppia $r=2$.

\par\noindent\hfill{\color{cyan!65!black}\scriptsize Pag. originale 18/521}\par\smallskip

% SIMULAZIONE / QUESITO 1 (pagina 18)
\begin{quiz}
Lo sviluppo di Mac Laurin, di ordine 2, della funzione $f(x)=\dfrac{1+\sin x}{\cos x}$ \`e:
\begin{enumerate}[label=\alph*),leftmargin=1.3cm,itemsep=3pt,topsep=5pt]
\item $f(x)=1+x+\frac12x^2$
\item $f(x)=2-x^2+o(x^2)$
\item $f(x)=1+x+\frac12x^2+o(x^2)$
\item $f(x)=1+x-\frac12x^2+o(x^2)$
\item $f(x)=1+x-\frac12x^2$
\end{enumerate}
\end{quiz}
% SOLUZIONE: c
% SPIEGAZIONE: $1/\cos x=1+x^2/2+o(x^2)$ e $\sin x=x+o(x^2)$.

% SIMULAZIONE / QUESITO 2 (pagina 18)
\begin{quiz}
Sia $f$ continua e derivabile nell\'intervallo $[4,6]$, con $f(4)=6$, $f(6)=0$. Allora necessariamente
\begin{enumerate}[label=\alph*),leftmargin=1.3cm,itemsep=3pt,topsep=5pt]
\item $\exists c\in(4,6):f(c)=0$
\item $\exists c\in(4,6):f\prime(c)=-3$
\item $\exists c\in[4,6]:f(c)=e-\pi$
\item $\exists c\in(4,6):f\prime(c)=3$
\item $\nexists c\in(4,6):f(c)>6$
\end{enumerate}
\end{quiz}
% SOLUZIONE: b
% SPIEGAZIONE: Per Lagrange esiste $c$ con $f\prime(c)=(0-6)/(6-4)=-3$.

% SIMULAZIONE / QUESITO 3 (pagina 18)
\begin{quiz}
Lo sviluppo di Mac Laurin, di ordine 4, della funzione $f(x)=\dfrac{3x^2}{1+3x}-3x^2-3x^4$ \`e:
\begin{enumerate}[label=\alph*),leftmargin=1.3cm,itemsep=3pt,topsep=5pt]
\item $-3x^4+o(x^4)$
\item $-9x^3-3x^4+o(x^4)$
\item $-9x^3+24x^4$
\item $-9x^3+24x^4+o(x^4)$
\item $-3x^3+8x^4+o(x^4)$
\end{enumerate}
\end{quiz}
% SOLUZIONE: d
% SPIEGAZIONE: $3x^2(1-3x+9x^2)-3x^2-3x^4=-9x^3+24x^4+o(x^4)$.

% SIMULAZIONE / QUESITO 4 (pagina 18)
\begin{quiz}
Sia $f$ continua e strettamente monotona, con $f(2)=7$, $f\prime(2)=1/7$, $f\prime(7)=1/2$. Allora $(f^{-1})\prime(7)=$
\begin{enumerate}[label=\alph*),leftmargin=1.3cm,itemsep=3pt,topsep=5pt]
\item $2$
\item $7$
\item non \`e possibile saperlo
\item $1/7$
\item $1/2$
\end{enumerate}
\end{quiz}
% SOLUZIONE: b
% SPIEGAZIONE: $(f^{-1})\prime(f(2))=1/f\prime(2)=7$.

% SIMULAZIONE / QUESITO 5 (pagina 18)
\begin{quiz}
Dallo sviluppo di Mac Laurin della funzione $f(x)=2+\log(5-\sin^2x)$ si pu\`o dedurre che $f$:
\begin{enumerate}[label=\alph*),leftmargin=1.3cm,itemsep=3pt,topsep=5pt]
\item ha un punto di massimo in $x=0$
\item punto di minimo in $x=0$
\item ha un flesso in $x=0$
\item \`e decrescente nell'intorno di $x=0$
\item \`e crescente nell'intorno di $x=0$
\end{enumerate}
\end{quiz}
% SOLUZIONE: a
% SPIEGAZIONE: $f(x)=2+\log5-x^2/5+o(x^2)$: massimo locale stretto.

\par\noindent\hfill{\color{cyan!65!black}\scriptsize Pag. originale 19/521}\par\smallskip

% SIMULAZIONE / QUESITO 6 (pagina 19)
\begin{quiz}
Lo sviluppo di Mac Laurin, di ordine 3, della funzione $\sinh(\log(x+1))$ \`e:
\begin{enumerate}[label=\alph*),leftmargin=1.3cm,itemsep=3pt,topsep=5pt]
\item $1+x-\frac12x^2+\frac{x^3}{3!}+o(x^3)+\frac1{3!}(1+x-\frac12x^2+\frac{x^3}{3!}+o(x^3))^3$
\item $x+\frac12x^2+\frac12x^3+o(x^3)$
\item $x-\frac12x^2-\frac12x^3+o(x^3)$
\item $1+x-\frac12x^2+\frac{x^3}{3!}+o(x^3)-\frac1{3!}(1+x-\frac12x^2+\frac{x^3}{3!}+o(x^3))^3$
\item $x-\frac12x^2+\frac12x^3+o(x^3)$
\end{enumerate}
\end{quiz}
% SOLUZIONE: e
% SPIEGAZIONE: $\sinh(\log(1+x))=((1+x)-(1+x)^{-1})/2=x-x^2/2+x^3/2+o(x^3)$.

% SIMULAZIONE / QUESITO 7 (pagina 19)
\begin{quiz}
Sia $w=\dfrac{5+i}{3-i}$. Allora $\operatorname{Re}(z)=$
\begin{enumerate}[label=\alph*),leftmargin=1.3cm,itemsep=3pt,topsep=5pt]
\item $5/3$
\item $-1$
\item $14$
\item $4/5$
\item $7/5$
\end{enumerate}
\end{quiz}
% SOLUZIONE: e
% SPIEGAZIONE: $w=(5+i)(3+i)/10=(14+8i)/10$, parte reale $7/5$. Nel testo \`e scritto $\operatorname{Re}(z)$ anzich\'e $\operatorname{Re}(w)$.

% SIMULAZIONE / QUESITO 8 (pagina 19)
\begin{quiz}
La funzione $f(x)=x^5+3e^x-2$ ha
\begin{enumerate}[label=\alph*),leftmargin=1.3cm,itemsep=3pt,topsep=5pt]
\item un solo zero positivo
\item non ha zeri
\item ha due zeri
\item ha almeno tre zeri
\item un solo zero negativo
\end{enumerate}
\end{quiz}
% SOLUZIONE: e
% SPIEGAZIONE: $f\prime(x)=5x^4+3e^x>0$; $f(0)=1>0$ e $f(x)\to-\infty$ a sinistra.

% SIMULAZIONE / QUESITO 9 (pagina 19)
\begin{quiz}
Sia $f(x)=-3+2(x-7)+5(x-7)^5+o((x-7)^5)$ lo sviluppo di Mac Laurin della funzione $f$. Quale delle seguenti propriet\`a \`e vera?
\begin{enumerate}[label=\alph*),leftmargin=1.3cm,itemsep=3pt,topsep=5pt]
\item la funzione \`e infinitesima, $x\to7$
\item la funzione \`e decrescente in un intorno di $x=7$
\item la retta tangente al grafico in $x=7$ \`e $y=-3+2(x-7)$
\item il punto $x=7$ \`e un punto critico
\item $f^{(5)}(7)=5/5!$
\end{enumerate}
\end{quiz}
% SOLUZIONE: c
% SPIEGAZIONE: $f(7)=-3$ e $f\prime(7)=2$, quindi la tangente \`e $y=-3+2(x-7)$.

% SIMULAZIONE / QUESITO 10 (pagina 19)
\begin{quiz}
Lo sviluppo di Mac Laurin di ordine 3 della funzione $\log(\cosh x)$ \`e:
\begin{enumerate}[label=\alph*),leftmargin=1.3cm,itemsep=3pt,topsep=5pt]
\item $\frac12x^2+o(x^3)$
\item $-\frac12x^2+o(x^3)$
\item $1+\frac12x^2+o(x^3)$
\item $1-\frac12x^2+o(x^3)$
\item $\frac12x^2$
\end{enumerate}
\end{quiz}
% SOLUZIONE: a
% SPIEGAZIONE: $\cosh x=1+x^2/2+O(x^4)$, dunque $\log\cosh x=x^2/2+O(x^4)$.

\par\noindent\hfill{\color{cyan!65!black}\scriptsize Pag. originale 20/521}\par\smallskip

% SIMULAZIONE / QUESITO 11 (pagina 20)
\begin{quiz}
Sia $w=\dfrac{5+i}{3-i}$. Allora il modulo di $w$ \`e
\begin{enumerate}[label=\alph*),leftmargin=1.3cm,itemsep=3pt,topsep=5pt]
\item $\sqrt3$
\item $\frac{\sqrt{65}}5$
\item $\sqrt{\frac{25}9+(-1)^2}$
\item $\sqrt{5^2+1^2}$
\item nessuna delle altre risposte
\end{enumerate}
\end{quiz}
% SOLUZIONE: b
% SPIEGAZIONE: $|w|=\sqrt{26}/\sqrt{10}=\sqrt{65}/5$.

% SIMULAZIONE / QUESITO 12 (pagina 20)
\begin{quiz}
Lo sviluppo di Mac Laurin, di ordine 3, della funzione $\log(\sinh x+1)$ \`e
\begin{enumerate}[label=\alph*),leftmargin=1.3cm,itemsep=3pt,topsep=5pt]
\item $x-\frac12x^2+\frac12x^3+o(x^3)$
\item $x-\frac12x^2+\frac13x^3+o(x^3)$
\item $x+\frac12x^2+\frac12x^3+o(x^3)$
\item $1+x-\frac12x^2+\frac12x^3+o(x^3)$
\item $x-\frac12x^2-\frac16x^3+o(x^3)$
\end{enumerate}
\end{quiz}
% SOLUZIONE: a
% SPIEGAZIONE: $\sinh x=x+x^3/6+o(x^3)$; in $\log(1+u)=u-u^2/2+u^3/3+\cdots$ il coefficiente cubico \`e $1/6+1/3=1/2$.

% SIMULAZIONE / QUESITO 13 (pagina 20)
\begin{quiz}
Il polinomio di Mac Laurin, di ordine 2, della funzione $f(x)=\dfrac{2e^{-x}-x^2}{1+x^2}$ \`e
\begin{enumerate}[label=\alph*),leftmargin=1.3cm,itemsep=3pt,topsep=5pt]
\item $2-2x-2x^2+o(x^2)$
\item $1-x-x^2$
\item $2+2x-4x^2$
\item $2-2x-2x^2$
\item nessuno degli sviluppi precedenti
\end{enumerate}
\end{quiz}
% SOLUZIONE: d
% SPIEGAZIONE: $2e^{-x}-x^2=2-2x+O(x^3)$ e $(1+x^2)^{-1}=1-x^2+o(x^2)$.

% SIMULAZIONE / QUESITO 14 (pagina 20)
\begin{quiz}
Quale delle seguenti propriet\`a NON \`e soddisfatta da $f\in C^{(7)}$ se $f(x)=5-4(x-2)^6+3(x-2)^7+o((x-2)^7)$ \`e lo sviluppo di Taylor in $x=2$?
\begin{enumerate}[label=\alph*),leftmargin=1.3cm,itemsep=3pt,topsep=5pt]
\item $f(x)$ \`e positiva nell'intorno di $x=2$
\item $x=2$ \`e punto critico
\item $f^{(7)}(2)=3\cdot7!$
\item $f$ \`e positiva e concava nell'intorno di $x=2$
\item $x=2$ \`e un flesso crescente
\end{enumerate}
\end{quiz}
% SOLUZIONE: e
% SPIEGAZIONE: Il termine dominante $-4(x-2)^6$ d\`a un massimo piatto in $2$, non un flesso crescente.

% SIMULAZIONE / QUESITO 15 (pagina 20)
\begin{quiz}
I primi tre termini del polinomio di Mac Laurin della funzione $f(x)=\dfrac{1+2x^2}{1-x^2}$ sono:
\begin{enumerate}[label=\alph*),leftmargin=1.3cm,itemsep=3pt,topsep=5pt]
\item $x+3x^2+6x^4$
\item $1-x^2+2x^4$
\item $1+3x^2+3x^4$
\item $1+x^2-x^4$
\item $x-x^3+x^5$
\end{enumerate}
\end{quiz}
% SOLUZIONE: c
% SPIEGAZIONE: $(1+2x^2)(1+x^2+x^4+\cdots)=1+3x^2+3x^4+\cdots$.

% SIMULAZIONE / QUESITO 16 (pagina 20)
\begin{quiz}
Qual \`e l\'insieme degli $\alpha\in\mathbb R$ per cui $\displaystyle\lim_{x\to0^+}\frac{x^\alpha+\log(1+x^2)}{x\sin x}=0$?
\begin{enumerate}[label=\alph*),leftmargin=1.3cm,itemsep=3pt,topsep=5pt]
\item $\{\alpha<2\}$
\item $\{\alpha=2\}$
\item $\{\alpha>2\}$
\item $\{\alpha=0\}$
\item $\varnothing$
\end{enumerate}
\end{quiz}
% SOLUZIONE: e
% SPIEGAZIONE: $\log(1+x^2)\sim x^2$ e $x\sin x\sim x^2$; per $\alpha>2$ il limite \`e $1$, per $\alpha=2$ \`e $2$, altrimenti non \`e $0$.
Risposta completata
